\chapter{La chaîne, en une page}\label{ch:chaine}

La méthode est la même sur tous les fonds de ce livre. On la résume ici pour pouvoir
ensuite parler de ce qui change d'une source à l'autre. Le détail des outils est dans
\emph{Mathématiques outillées}.

% Chapitre 1 : la chaîne, de l'inventaire aux trouvailles.
\begin{figure}[tbp]
\centering
\begin{tikzpicture}[x=1cm,y=1cm, node distance=4mm]
  \node[boite, text width=2.1cm] (inv) at (0,0) {Inventaire\\[-1pt]\scriptsize les cotes du détenteur};
  \node[boite, text width=2.1cm, right=of inv] (copie) {Copie locale\\[-1pt]\scriptsize ignorée par \texttt{git}};
  \node[boite, text width=2.1cm, right=of copie] (tr) {Transcription\\[-1pt]\scriptsize lots de 20 pages};
  \node[boite, text width=2.1cm, below=9mm of tr] (rel) {Seconde passe\\[-1pt]\scriptsize puis une personne};
  \node[boite, text width=2.1cm, left=of rel] (mod) {Lecture modernisée\\[-1pt]\scriptsize un autre fichier};
  \node[boite, text width=2.1cm, left=of mod] (lean) {Lean, trouvailles\\[-1pt]\scriptsize contre la littérature};
  \draw[fleche] (inv) -- (copie);
  \draw[fleche] (copie) -- (tr);
  \draw[fleche] (tr) -- (rel);
  \draw[fleche] (rel) -- (mod);
  \draw[fleche] (mod) -- (lean);
  \node[boite, fill=white, dashed, text width=6.6cm, below=9mm of mod] (site)
    {Le site : la transcription à côté du fac-similé, servi par la source};
  \draw[fleche, dashed] (rel.south) |- (site.east);
  \draw[fleche, dashed] (mod) -- (site);
\end{tikzpicture}
\caption{La chaîne, la même sur tous les fonds. Les étapes du haut dépendent de la
source (format, serveur, débit) ; celles du bas dépendent des droits.}
\label{fig:chaine}
\end{figure}


\begin{enumerate}
  \item \textbf{L'inventaire.} Un script lit le catalogue de l'institution (ses cotes, ses
        dates, son nombre de pages) et en fait un fichier de données du dépôt. Les cotes,
        les datations et les descriptions sont celles du détenteur, citées comme telles ; un
        regroupement qu'on ajoute (un « cahier », une période) est signalé comme le sien.
  \item \textbf{La copie de travail locale.} Pour lire une page de près, l'assistant a
        besoin de l'image en pleine résolution. On en tire une copie locale, ignorée par
        \texttt{git}, lot par lot, avec des pauses entre les requêtes. Elle sert à
        transcrire et n'est jamais publiée.
  \item \textbf{La transcription}, par lots de vingt pages (quarante pour des
        dactylographies), en LaTeX, avec un appareil critique : ce qui est lu, ce qui est
        deviné (\verb|\uncertain{}|), ce qui est illisible (\verb|\ill{}|), ce qui est
        barré, ajouté, en marge. Le texte et l'appareil sont deux registres séparés.
        Chaque fichier dit dans son en-tête qui l'a lu et quand.
  \item \textbf{La seconde passe}, sur des agrandissements de la page, qui reprend les
        lectures douteuses. Elle reste une relecture par la machine : la vérification par
        une personne, contre le fac-similé, est déclarée à part et lot par lot.
  \item \textbf{La lecture modernisée}, un fichier distinct : un résumé, puis les
        mathématiques en notation et en vocabulaire actuels, chaque traduction de notation
        signalée en note.
  \item \textbf{Le site}, où la transcription s'affiche à côté du fac-similé, servi par
        l'institution elle-même chaque fois que c'est possible.
  \item \textbf{Lean}, pour les énoncés précis de la lecture modernisée qu'on peut
        démontrer, et \textbf{les trouvailles}, chacune écrite comme une affirmation sur
        la littérature (« introuvable dans tel et tel ouvrage »), jamais comme une
        priorité.
\end{enumerate}

Chaque étape a son état, lu dans les fichiers et affiché sur le site : à faire, en cours,
transcrit, relu par la machine, vérifié par une personne. La dernière case est presque vide
sur tous les fonds de ce livre, et le site le dit.

\section{D'une source à l'autre}

Rien dans cette liste ne dépend du sujet mathématique. Presque tout dépend en revanche de
la source : la copie de travail dépend du format des images et de ce que le serveur
tolère ; l'affichage à côté du fac-similé dépend des en-têtes que le serveur envoie ; la
publication des transcriptions dépend des droits ; la transcription elle-même dépend de la
main, de la langue, de la numérotation des pages. Le chapitre suivant fait de ces questions
une grille. Les chapitres de la deuxième partie l'appliquent, source par source.
